\documentclass[12pt]{article}

\begin{document}
\textbf{C02. Topics in Nonlinear Dynamics}

\hfill

\textbf{2.1 The Cardiovascular System}

\hfill


Applications of nonlinear dynamics to the cardiovascular system are typically
related to the electrophysiology of the heart such that the phenomenon of interest is usually events that are of the time intervals less than 1 min. Of particular interest is
the heart rate variation and arrhythmias. Nonlinear methods provide a different perspective
into electrophysiology and are typically not taught in the classical medical
physiology textbook. In particular, physiology can explain and identify the origin of
some arrhythmias, but the most common, the sinus arrhythmia, remains to be
explained. The cardiac rhythm was generally introduced in Chap. 1 as the concept
of the bio-oscillator. A nonlinear oscillator may be used to analyze many known
cardiac rhythms normal and pathological. Moreover, since the cardiac pacemaker is
a repetitive process, the concept of the Fractal can be employed. Cardiac electrophysiology
has enjoyed a long history of research, so that information is available
at the heart level, the cellular level, and the molecular level. Nonlinear dynamics
will be employed in this book to revisit each structural level with a new vantage
point. A focal point of debate is the origin of cardiac fibrillation. Briefly, fibrillation
is known as an uncoordinated electrical and mechanical activation of the cardiac
muscle that leads to death due to net zero pumping of blood. Some dynamicists have
theorized that fibrillation is Chaos. These discussions will be covered in more detail
in later chapters after the reader becomes more familiar with nonlinear dynamics.
Continuing along with our survey of nonlinear dynamics in physiology, move next
to the pulmonary system.

\hfill

\textbf{2.2 Respiratory Physiology}

\hfill

Recall from Chap. 1 that the Fractal concept applies well to the modeling of the geometry
of the respiratory system. Also, the control of breathing lends itself well to the
application of nonlinear dynamic methods. Of course, everyone is familiar with
breathing being a periodic process. It would be simple to suggest that breathing is
stimulated via pacemaker cells just as in the heart. Surprisingly, the respiratory system
possesses no such respiratory oscillator. Instead, the use of a nonlinear feedback control
system explains the source of oscillations (Glass \& Mackey, 1988). Physiologists
are familiar with the homeostatic reflex regulation of breathing. Reflex regulation is a
negative feedback control system that is employed extensively in human physiology.
By adding nonlinear affecters and feedback delay, it is possible to cause the respiratory
regulation system to become unstable and oscillate. This kind of treatment of the
respiratory system explains the source of oscillatory breathing. This approach further
explains some pathology of breathing such as respiratory arrest and Cheyne’s stokes
phenomenon and other abnormal breathing patterns.

\hfill

\textbf{2.3 Population Dynamics}

\hfill


Probably one of the first applications of nonlinear dynamics to biology is that of the
study of population dynamics, that is, to analyze how various animal populations
vary in time. Of recent interest is the study of how diseases may spread through a
population. The analysis of disease spread is a subject of epidemiology population
dynamics analysis applied well in the case of a confined population. For example, a
human population that is constrained within the borders of a country. Many readers
are familiar with epidemiologist presentations of the spread of COVID infections
during 2020 (Giordano et al., 2020). The theory behind the infection–time curves
rests on the analysis of population dynamics. Additionally, population dynamics
may also be applied to populations of cells. For example, the dynamics of the human
blood cell population. Human cell dynamics is useful to help analyze various types
of anemia in white and red cell concentration. The analytical knowledge of these
cell populations can be used for more effective treatment of disease that causes
extreme low-cell concentrations.

\hfill

In summary, nonlinear dynamics analysis applies well to much of human physiology.
We will explore more of these applications as we move through this book. At
this time, it is useful to examine some of the analytical methods that are valuable in
our study such as basic population models.

\hfill

2.4 Population Model Solution

\hfill

To illustrate some of the fundamental methods by which the nonlinear biological
system may be analyzed, a simple population model is now introduced and solved.
A characteristic feature of unrestrained population growth is that the rate of growth
is proportionate to the magnitude of the population itself. Mathematically, this can
be written as

\[
P'(t) = aP(t)
\]

where P(t) is the population versus time, and P' is its time derivative; “a” is a constant
of proportionality that must be calibrated to the population. The solution of
this differential equation is easily determined to be

\[
P = PP_0 e^{at}
\]

The constant $P_0$ is the initial population. Following this model, it is found that the
population is unstable, that is, as t increases, the population increases exponentially
to infinity for $a > 0$. Reexamining the model for the case where $a < 0$, it is found that
the model would predict that population decreases exponentially to zero as t
increases. Hence, for either sign of the model, results are unreasonable. A sample
population growth curve is shown in Fig. 2.1 for a positive value of a, where infinite
population growth occurs.

\hfill

This model is flawed in the final value prediction of the model. There are only
two final values: infinite population or a zero population. The model is flawed in
these outcomes, but the model predictions are not unreasonable over short time
intervals. So, rather than abandoning the model completely, look at a modification.
Returning to Eq. 2.1 of the model, it is a linear model that permits the population to
grow in direct proportion to its size. This problem can be repaired by adding a term
that prevents growth when the population becomes too large. So, modify the population
model now to limit growth according to a new differential equation,

\[
\frac{dP}{dt} = a P (L-P)
\]

where L is the population limit. Here, we will use a normal value of L = 1. For
Eq. 2.3, it can be seen that as the population grows larger than L, the derivative of
population growth will turn negative and prevent further growth. Or, when L = P. dP/dt = 0, so that the steady state P of this new model is for population to converge to
L. Let’s confirm this by solving Eq. 2.3 for P(t). At this point, we run into a new
problem. That is, after expanding the right side of Eq. 2.3 we see that a term $aP^2$ emerges to result in a nonlinear differential equation. Hence, the reader now sees
our first example of a nonlinear dynamic system in this book. There is much to learn
from this very simple nonlinear model. To solve this equation, Euler method is now
introduced. This will be a computational solution of the differential equation.

\hfill

Start by defining two adjacent population points in time: population at the current time is $P(t_i)$, then, after a small change in time $\Delta t. P_{i+1} = P(t_i + \Delta t)$. Assuming
that the time change can be made infinitely small, the derivative of population then
can be written as

\[
\frac{dP}{dt} = \frac{P_{i+1} - P_i}{\Delta t}
\]


Using the definitions and substituting them into Eq. 2.3, the new model is converted
to Euler computational form to be Eq. 2.4:

\[
P_{i+1} = \Delta t a P_i(1-P_i) + P_i
\]

This form of our new population model is known to those who study nonlinear
dynamics or population modeling as the logistic equation [REF]. Equation 2.4 is
now in the form such that we enter the current population value Pi and simply compute
the next value of population Pi+1 Hence, we now have an iterative solution for
the original differential equation model Eq. 2.4 A sample iteration of Eq. 2.4 is
provided in Fig. 2.2. The properties of this equation are very illustrative of nonlinear
dynamic systems. These features will be examined in more detail in a later chapter,
but for now notice that the new population model is functioning as required. In
Fig. 2.2, the vertical axis represents population, and the horizontal axis is time,
where each value n indicates a time step of $\Delta t$. Observe that population now converges
to a specific level of about 0.5. As you ‘ll find later, no matter what the initial
population starts at, this model will converge to the .5 level. This was determined by
the choice of parameter a = 2.

\hfill

As you study this simple nonlinear model, it is worth noting that the computational
form of the model is very simple when it is presented as an iteration as above.
Consider that the model constants may be combined, and that the resulting solution
is P(t+Δt) = αP (t) with all constants combined to be α or as Pi + 1 = αPi. It is clear
that the nonlinear system solution has been reduced to a very elementary form. This
approach is what biology has been optimized to perform very well where one very
simple process is then performed repeatedly to yield a much more complex function.
This simplification procedure will be highlighted throughout this book. Thus,
the student of nonlinear dynamics need not learn a diverse set of procedures to study
the nonlinear dynamic biological systems. The analytical tools that you will learn
here will have general application. We have seen that an important biodynamic process
is the oscillator. Or, otherwise known as the pacemaker. Hence, it is useful to
introduce bio-oscillations into this chapter as a special nonlinear process.
2.5 Bio-oscillation

\hfill

The vertical axes represent population value and the horizontal axes represent
time. In this graph, a value of a = 2 and initial population of 0.6 were used.

\end{document}